\akhead{Topic 7 — Nonlinear Equations and Systems of Equations}

\ans{1}{1}{$8$}{Taking square roots gives $x=\pm 8$; the condition $x>0$
leaves $x=8$.}
\ans{1}{2}{A}{A product is $0$ when a factor is $0$, so $x=3$ or $x=-7$. The
smaller is $-7$.}
\ans{1}{3}{D}{Two numbers with product $12$ and sum $7$ are $3$ and $4$, so
$(x-3)(x-4)=0$ and the solutions are $3$ and $4$.}
\ans{1}{4}{$9$}{Take square roots: $x-4=\pm 5$, so $x=9$ or $x=-1$. Only
$x=9$ is positive.}
\ans{1}{5}{B}{With $a=1$, $b=4$, $c=3$, $b^2-4ac=16-12=4$.}
\ans{1}{6}{A}{The discriminant is $2^2-4(1)(5)=-16$, which is negative, so the
equation has no real solutions.}
\ans{1}{7}{$9$}{Square both sides: $x+7=16$, so $x=9$. Check:
$\sqrt{16}=4$.}
\ans{1}{8}{C}{$x-5=3$ gives $x=8$ and $x-5=-3$ gives $x=2$. The greater is
$8$.}
\ans{1}{9}{D}{Factor: $(2x+1)(x-3)=0$, so $x=-\tfrac12$ or $x=3$. The greater
is $3$.}
\ans{1}{10}{$5$}{Cross multiply: $x+4=3(x-2)=3x-6$, so $2x=10$ and $x=5$.
Neither denominator is $0$ at $x=5$.}
\ans{1}{11}{B}{Move the constant: $x^2+8x=-3$. Add $(8/2)^2=16$ to both sides:
$(x+4)^2=13$, so $b=13$.}
\ans{1}{12}{C}{The sum of the solutions of $ax^2+bx+c=0$ is $-b/a$, here
$-(-6)/1=6$. (The solutions are $2$ and $4$.)}
\ans{1}{13}{C}{Squaring gives $2x+3=x^2$, so $x^2-2x-3=0$ and $x=3$ or
$x=-1$. Substituting $x=-1$ gives $\sqrt{1}=1\neq -1$, so $-1$ is extraneous
and the solution is $3$.}
\ans{1}{14}{$2$}{Multiply by $x-1$: $x+3=5x-5$, so $4x=8$ and $x=2$. Since
$x\neq 1$, $x=2$ is valid.}
\ans{1}{15}{B}{With width $w$, $w(w+4)=96$, so $w^2+4w-96=0$ and
$(w+12)(w-8)=0$. A width must be positive, so $w=8$ meters.}
\ans{1}{16}{C}{Set $-16t^2+48t=0$, so $-16t(t-3)=0$ and $t=0$ or $t=3$. The
kick is at $t=0$, so the ball lands at $t=3$ seconds.}
\ans{1}{17}{D}{Substitute: $x^2-2x-3=x-3$, so $x^2-3x=0$ and $x(x-3)=0$. The
$x$-coordinates are $0$ and $3$; the greater is $3$.}
\ans{1}{18}{$10$}{Exactly one real solution means $b^2-4ac=0$, so
$k^2-4(1)(25)=0$ and $k^2=100$. Since $k>0$, $k=10$.}
\ans{1}{19}{B}{If the greater integer is $n$, then $n(n-2)=168$, so
$n^2-2n-168=0$ and $(n-14)(n+12)=0$. The integers are positive, so $n=14$.}
\ans{1}{20}{C}{Substitute $y=x+1$: $x^2+(x+1)^2=25$, so $2x^2+2x-24=0$ and
$x^2+x-12=0$. Then $(x-3)(x+4)=0$, giving $x=3$ or $x=-4$; the greater is
$3$.}
\ans{1}{21}{C}{Setting the expressions equal gives $x^2+4x+k=2x$, so
$x^2+2x+k=0$. One intersection point means $2^2-4(1)(k)=0$, so $k=1$.}
\ans{1}{22}{$3$}{Either $x^2-4=5$, giving $x^2=9$ and $x=\pm 3$, or
$x^2-4=-5$, giving $x^2=-1$, which has no real solution. The positive
solution is $3$.}

\ans{2}{1}{C}{Factor: $(x+6)(x-2)=0$, so $x=-6$ or $x=2$. The positive
solution is $2$.}
\ans{2}{2}{$7$}{Square both sides: $4x-3=25$, so $4x=28$ and $x=7$. Check:
$\sqrt{25}=5$.}
\ans{2}{3}{B}{$3x+2=11$ gives $x=3$; $3x+2=-11$ gives $3x=-13$, so
$x=-\tfrac{13}{3}$.}
\ans{2}{4}{C}{Cross multiply: $5(x-2)=2(x+1)$, so $5x-10=2x+2$ and $3x=12$,
giving $x=4$. Neither denominator is $0$ there.}
\ans{2}{5}{D}{Rewrite as $x^2-5x-24=0$ and factor: $(x-8)(x+3)=0$, so $x=8$
or $x=-3$. The positive solution is $8$.}
\ans{2}{6}{B}{The discriminant is $(-4)^2-4(1)(4)=0$, so there is exactly one
real solution, namely the double root $x=2$.}
\ans{2}{7}{B}{Substitute $y=3$: $x^2-4x+7=3$, so $x^2-4x+4=0$ and
$(x-2)^2=0$. The only $x$-coordinate is $2$.}
\ans{2}{8}{$12$}{$x^2-10x=-18$; add $(10/2)^2=25$ to both sides to get
$(x-5)^2=7$. So $a=5$, $b=7$, and $a+b=12$.}
\ans{2}{9}{C}{Set $-5t^2+20t+25=0$, divide by $-5$: $t^2-4t-5=0$, so
$(t-5)(t+1)=0$. Time cannot be negative, so $t=5$ seconds.}
\ans{2}{10}{B}{For $x\neq 3$, $\dfrac{x^2-9}{x-3}=x+3$, so $x+3=8$ and
$x=5$.}
\ans{2}{11}{D}{Solve $x^2-6x+5=12$, that is $x^2-6x-7=0$, so $(x-7)(x+1)=0$
and $x=7$ or $x=-1$. The greater is $7$.}
\ans{2}{12}{$4$}{Substitute $y=2x$: $x^2+4x^2=20$, so $x^2=4$ and $x=\pm 2$.
The solutions are $(2,4)$ and $(-2,-4)$, so the greatest $y$ is $4$.}
\ans{2}{13}{B}{Factor: $(3x+1)(x-2)=0$, so $x=-\tfrac13$ or $x=2$. Of the
options, only $-\tfrac13$ is a solution.}
\ans{2}{14}{A}{Half of $-12$ is $-6$ and $(-6)^2=36$, so
$x^2-12x+40=(x^2-12x+36)+4=(x-6)^2+4$.}
\ans{2}{15}{$-7$}{Substitute $x=3$: $9+3k+12=0$, so $3k=-21$ and $k=-7$.}
\ans{2}{16}{C}{Exactly one real solution means $b^2-4(4)(9)=0$, so
$b^2=144$. Since $b$ is positive, $b=12$.}
\ans{2}{17}{C}{Setting the expressions equal gives $x^2+6x+c=2x+1$, so
$x^2+4x+(c-1)=0$. One intersection means $4^2-4(c-1)=0$, so $c-1=4$ and
$c=5$.}
\ans{2}{18}{$-1$}{Squaring gives $2x+11=x^2+8x+16$, so $x^2+6x+5=0$ and
$x=-1$ or $x=-5$. At $x=-5$ the right side is $-1$ while the radical is $1$,
so $-5$ is extraneous; the solution is $-1$.}
\ans{2}{19}{A}{Two distinct real solutions require $10^2-4k>0$, that is
$k<25$. Of the options only $24$ satisfies this.}
\ans{2}{20}{B}{Substituting $y=k$ gives $x^2=9-k^2$, which has exactly one
solution when $9-k^2=0$. Since $k>0$, $k=3$: the line touches the circle at
its top point $(0,3)$.}
\ans{2}{21}{$12$}{The two side lengths sum to $34/2=17$ and multiply to $60$,
so they are the solutions of $t^2-17t+60=0$, that is $(t-5)(t-12)=0$. The
longer side is $12$ centimeters.}
\ans{2}{22}{$\dfrac{2}{3}$}{Either $x-3=2x+1$, giving $x=-4$, or
$x-3=-(2x+1)$, giving $3x=2$ and $x=\tfrac23$. The positive solution is
$\tfrac23$.}
